Para garantizar la reproducción de los resultados, los experimentos se realizaron en un sistema con las siguientes especificaciones:

\subsection*{Especificaciones}
\begin{itemize}
  \item \textbf{Equipo:} \texttt{LENOVO ThinkPad X13 Gen 3}.
  \item \textbf{CPU:} \texttt{Intel Core i7-1265U (4.8 GHz)}.
  \item \textbf{RAM:} \texttt{16 GB}.
  \item \textbf{Disco:} \texttt{SSD NVMe 512 GB}.
  \item \textbf{SO:} \texttt{Fedora Linux 43 (Workstation Edition)} con \texttt{Linux 6.19.10-200.fc43.x86\_64}.
  \item \textbf{Compilador:} \texttt{g++ 11.4.0} con flags \texttt{-std=c++17 -O2 -g -Wall -Wextra -pedantic}.
\end{itemize}

La compilación y ejecución se automatizó con \texttt{Makefile}. Los tiempos se midieron en nanosegundos $(ns)$ con \texttt{<chrono>} \parencite{cppreference_chrono} y el pico de memoria en kilobytes (KB) mediante \texttt{getrusage} en Linux. \parencite{linux_getrusage}

